\BeleskeID{08-s049}
% element: 08-s049:e03
% element: 08-s049:notes
Izraz za izlazni napon teži ka konačnoj vrednosti $U$, ali je dostiže tek u granici $t\to\infty$. Zato tačno matematičko trajanje prelaznog procesa nije praktično za crtanje i procenu završetka.

Posle pet vremenskih konstanti preostalo odstupanje čini $e^{-5}$ početnog odstupanja, a ostvareni deo promene je $1-e^{-5}\approx0.9933$, odnosno približno 99.3\%. Pravilo $5\tau$ predstavlja praktičan kriterijum, a ne tačno iščezavanje eksponencijalnog člana.
